% !TeX root = ../../Book1.tex
% 纲要标题：### 9.1 自由阿贝尔群
% 纲要位置：计划纲要.md，第 358 行。

\section{自由阿贝尔群}
\label{b1:sec:0901}

% 纲要内容（仅作写作计划，尚非教材正文）：
%
% * 整数格与自由阿贝尔群
% * 子群仍自由及有限秩情形
% * 秩及其不变性
%

\subsection{整数格与自由阿贝尔群}
\label{b1:subsec:0901:01}

阿贝尔群的乘法可以改写为加法，生成元之间的关系因而成为整数线性关系。本章统一用 \(0\) 表示单位元，用 \(mx\) 表示元素 \(x\) 的整数倍。先从没有非平凡整数关系的群开始。

\begin{definition}\label{b1:def:free-abelian}
设 \(I\) 为集合，令
\[
\mathbb Z^{(I)}
=\{\,(n_i)_{i\in I}\mid n_i\in\mathbb Z,\ %
n_i\ne0\text{ 的指标只有有限个}\,\}.
\]
逐分量加法使它成为阿贝尔群。与某个 \(\mathbb Z^{(I)}\) 同构的群称为\term[ziyou-abeier-qun]{自由阿贝尔群}[free abelian group]。其中标准向量 \(e_i\) 的对应元素构成一组基：每个元素均可唯一写为有限和 \(\sum n_i e_i\)。
\end{definition}
\symidx{\(\mathbb Z^{(I)}\)：有限支集整数族}
有限的 \(I\) 给出熟悉的整数格 \(\mathbb Z^n\)。无限指标时必须限定有限支集；允许任意整数族所得的全直积是不同的构造，本章不把它当作自由阿贝尔群的定义。

\begin{proposition}\label{b1:prop:free-abelian-universal}
对任意阿贝尔群 \(A\) 和映射 \(f:I\to A\)，存在唯一同态
\(\widetilde f:\mathbb Z^{(I)}\to A\)，使 \(\widetilde f(e_i)=f(i)\)。
\end{proposition}
\begin{proof}
定义 \(\widetilde f(\sum n_i e_i)=\sum n_i f(i)\)。表示的唯一性保证良定义；和只有有限项，交换性保证加法相容。生成元的像已经指定，所以这样的同态唯一。
\end{proof}
这与自由群的泛性质相似，区别在于目标限为阿贝尔群。比如从 \(\mathbb Z^2\) 到群 \(G\) 指定两个基向量的像 \(x,y\)，必须要求 \(xy=yx\)；从 \(F_2\) 出发则没有这个要求。

\ProseParagraphBreak
在 \(\mathbb Z^2\) 中，\((1,0),(1,1)\) 构成一组基，因为
\((u,v)=(u-v)(1,0)+v(1,1)\)，系数唯一。而 \((2,0),(0,1)\) 虽然整数线性无关，却只生成第一坐标为偶数的子群。整数格的“无关”与“生成”必须分别检查，不能把实向量空间中的一组基自动当作整数格的基。

\subsection{自由阿贝尔群的子群}
\label{b1:subsec:0901:02}

\begin{theorem}\label{b1:thm:free-abelian-subgroup}
自由阿贝尔群的每个子群仍是自由阿贝尔群。若 \(H\le\mathbb Z^n\)，则 \(H\) 有一组至多 \(n\) 个元素的基。
\end{theorem}
\begin{proof}
先证明有限秩情形，对 \(n\) 归纳。\(n=0\) 显然。对 \(n\ge1\)，把 \(H\) 投影到最后一个坐标，其像是 \(\mathbb Z\) 的子群，所以为 \(d\mathbb Z\)，其中 \(d\ge0\)。核
\[
K=H\cap(\mathbb Z^{n-1}\times\{0\})
\]
由归纳假设自由且秩至多 \(n-1\)。若 \(d=0\)，则 \(H=K\)。若 \(d>0\)，选 \(h\in H\) 的末坐标为 \(d\)。任意 \(x\in H\) 的末坐标为 \(md\)，于是 \(x-mh\in K\)；而 \(mh\in K\) 迫使 \(md=0\)，进而 \(m=0\)。所以 \(H=K\oplus\mathbb Zh\)，在 \(K\) 的基后加上 \(h\) 即可。

一般情形在通常的选择公理下，把母群的一组基良序为 \((e_\alpha)_{\alpha<\kappa}\)。令 \(F_\alpha\) 由所有 \(e_\beta\)（\(\beta<\alpha\)）生成，令 \(H_\alpha=H\cap F_\alpha\)。从 \(H_{\alpha+1}\) 取 \(e_\alpha\) 的坐标，像为 \(d_\alpha\mathbb Z\)。若 \(d_\alpha>0\)，选 \(h_\alpha\) 的此坐标为 \(d_\alpha\)，完全相同的减去整数倍步骤给出
\[
H_{\alpha+1}=H_\alpha\oplus\mathbb Zh_\alpha.
\]
若 \(d_\alpha=0\)，两项相等。对极限指标 \(\lambda\)，有限支集保证
\(H_\lambda=\bigcup_{\alpha<\lambda}H_\alpha\)。

所选全部 \(h_\alpha\) 生成 \(H\)：逐级使用上式，并在极限步骤取并即可；任何给定元素只涉及有限项。它们也整数线性无关，因为任何有限关系中指标最大的 \(h_\alpha\) 在 \(e_\alpha\) 位置的系数为其整数系数乘 \(d_\alpha\)，更早的项在该位置全为零，所以该系数必须为零；依次向前得到所有系数为零。这便是一组自由基。
\end{proof}
证明并未断言子群的基总能扩充为母群的基。子群 \(2\mathbb Z\le\mathbb Z\) 的基为 \(2\)，却不能成为 \(\mathbb Z\) 的基的一部分；商群在这里出现了二阶挠元。作为可跟算的二维例子，
\[
H=\{\,(u,v)\in\mathbb Z^2\mid u+v\equiv0\pmod3\,\}
\]
在第二坐标的投影为 \(\mathbb Z\)，可选 \(h=(-1,1)\)；核为 \(3\mathbb Z\times\{0\}\)。所以 \((3,0),(-1,1)\) 是 \(H\) 的一组基。

\subsection{秩及其不变性}
\label{b1:subsec:0901:03}

基的个数必须与选择无关，才可用来比较群。

\begin{proposition}\label{b1:prop:free-abelian-rank}
自由阿贝尔群任意两组基的基数相同，称为\term[ziyou-abeier-qun-de-zhi]{自由阿贝尔群的秩}[rank of a free abelian group]。
\end{proposition}
\begin{proof}
若一组基有有限个 \(n\) 元，则每个元素模 \(2A=\{2x\mid x\in A\}\) 后，可唯一把各坐标减为 \(0\) 或 \(1\)，所以 \(|A/2A|=2^n\)。另一组有限基因而具有相同大小。若 \(A\) 有有限生成集，则这些生成元相对于任意一组基只涉及有限多个基元素；其余基元素不可能由它们生成，所以每组基都有限。

对两组无限基 \(B,C\)，把 \(C\) 中的每个元素写为 \(B\) 的有限线性组合。这些支集的并必须覆盖 \(B\)，故 \(|B|\le |C|\)；交换角色得到反向不等式。这与自由群的无限秩证明使用相同的有限支集和基数比较。
\end{proof}
整数线性无关最多只能给出秩的下界。例如前面的指数三子群与 \(\mathbb Z^2\) 都有秩二，但并不相等。另一方面，秩为一的自由阿贝尔群必为无限循环群，因此加法群 \(\mathbb Q\) 若自由就不可能仅凭“嵌在一条实直线中”断言它秩为一；下一节会直接证明它根本不自由。

\ProseParagraphBreak
有限生成阿贝尔群都有从某个 \(\mathbb Z^n\) 出发的满同态。其核由本节定理仍是有限秩自由阿贝尔群，于是关系可以选为有限多个整数向量。这是后面使用有限整数矩阵进行分类的依据；有限生成的关系并非一开始就应当默认为已知。
